\documentclass{article}
\usepackage{../math_template}

\author{Jonathan Kasongo}
\title{International Math Comp 2002 Day 2 --- P2/5}

\begin{document}
\maketitle

\begin{problem}{International Math Comp 2002 Day 2 --- P2/5}
Two hundred students participated in a mathematical contest. They had 6
problems to solve. It is known that each problem was correctly solved by at
least 120 participants. Prove that there must be two participants such that
every problem was solved by at least one of these two students.
\end{problem}

\begin{solution}{Solution}
\textit{Claim: The student(s) who solved the most problems, solved
at least 4 problems.}\vspace{0.2cm}

\textit{Proof:} Note that there were at least $120\times 6 = 720$ distinct
solutions submitted during the contest. If the student who solved the most
problems solved $s < 4$ problems, then there had to have been at most
$200\times 3 = 600$ distinct solutions submitted, a contradiction. $\square$
\vspace{0.2cm}

Suppose WLOG that the student who solved the most problems, correctly solved
problems 1,2,3 and 4. We wish to show that there exists a student who
correctly solved problems 5 and 6. However due to the pigeonhole principle
at least $120\times 2 - 200 = 40$ students solved both problems 5 and 6.
Thus there does indeed exist two participants such that every problem
was solved by at least one of these two students. $\blacksquare$
\end{solution}

\end{document}
